
\usepackage{cmap}
\usepackage[T2A]{fontenc}
\usepackage[utf8]{inputenc}         % file encoding
\usepackage{substitutefont}         % so we can use fonts other than those specified in babel
\usepackage[english]{babel}         % default language for document
\usepackage[scaled=0.925]{XCharter} % Подключение русифицированных шрифтов XCharter
\usepackage[bitstream-charter]{mathdesign} % Согласование математических шрифтов

\usepackage{
    %amsfonts,
    mathtools,
    mathtext,
    cite,
    enumerate,
    float,
    textcomp
}                                   % some packages 
%\usepackage[font={small}]{caption}
\usepackage{hhline}                 % beautiful links
\usepackage{hyperref}               % beautiful links
\usepackage{pdflscape}              % landscape pages
\usepackage{longtable}              % longtable support
\usepackage{multirow}               % vertical cell in table
\usepackage{bigstrut}               % big strut
\usepackage{array}                  % additional cell aligh
\usepackage{indentfirst}            % first line indent
\usepackage{gensymb}                % symbols

%\usepackage[nottoc]{tocbibind} %  do we need bibliography in toc
%----------------------------------------------
%\usepackage{mhchem}        
%===============  Font for tables  ============
%\let\oldtabular\tabular
%\renewcommand{\tabular}{\small\oldtabular}

\hypersetup{colorlinks=true,
    linkcolor=blue,
    citecolor=blue, 
    filecolor=blue, 
    urlcolor=blue, 
    pdftitle={Cinelerra GG Infinity Manual}, 
    pdfauthor={Cinelerra Authors},
    pdfsubject={Video Editing}, 
    pdfkeywords={Cinelerra, Good Guy}
} % pdf properties
\usepackage[pdftex]{graphicx}       % do we need some figures in our pdf 
\graphicspath{{images/}}            % path to images
\usepackage{tikz}                   % drawing package
\usepackage{nameref}                % use \nameref{} to set reference to chapter neme.
%\usepackage{noto}


%----------------------------------------------------------------------
\usepackage{listings}               % include code 
\lstset{                            % begin settings
  %language=R,                      % the language of the code
  inputencoding=utf8,
  basicstyle=\ttfamily\footnotesize,         % the size of the fonts that are used for the code
  numbers=left,                     % where to put the line-numbers
  numberstyle=\tiny\color{black},   % the style that is used for the line-numbers
  stepnumber=1,                     % the step between two line-numbers. If it's 1, each line
                                    % will be numbered
  numbersep=5pt,                    % how far the line-numbers are from the code
  %backgroundcolor=\color{white},   % choose the background color. You must add \usepackage{color}
  showspaces=false,                 % show spaces adding particular underscores
  showstringspaces=false,           % underline spaces within strings
  showtabs=false,                   % show tabs within strings adding particular underscores
  frame=lines,                      % adds a frame around the code
  %frame=single,                    % adds a frame around the code
  rulecolor=\color{black},          % if not set, the frame-color may be changed on line-breaks within not-black text (e.g. commens (green here))
  tabsize=2,                        % sets default tabsize to 2 spaces
  captionpos=b,                     % sets the caption-position to bottom
  breaklines=true,                  % sets automatic line breaking
  breakatwhitespace=false,          % sets if automatic breaks should only happen at whitespace
  title=\lstname,                   % show the filename of files included with \lstinputlisting;
                                    % also try caption instead of title
  keywordstyle=\color{blue},        % keyword style
  commentstyle=\color{gray},        % comment style
  stringstyle=\color{black},        % string literal style
  %backgroundcolor=\color{green!10},
  escapeinside={\%*}{*)},           % if you want to add a comment within your code
  extendedchars=\true,
  %keepspaces = true                %!!!! spaces in comments
  texcl=true,
  morekeywords={*,...}              % if you want to add more keywords to the set
}
%---------------------------------------------------------------------------
\makeatletter
\renewcommand{\@biblabel}[1]{#1.}
%---------------------------------------------------------------
%
\usepackage{tikz}                   % Powerfull ploting language
\usetikzlibrary{                    % Libraries for TiKz
    positioning,
    arrows,
    shapes,
    shadows
} 



% Title page definition
%\providecommand{\HUGE}{\Huge}% if not using memoir
\newlength{\drop}% for my convenience
%% specify the Webomints family
%\newcommand*{\wb}[1]{\fontsize{#1}{#2}\usefont{U}{webo}{xl}{n}}
%% select a (FontSite) font by its font family ID
\newcommand*{\FSfont}[1]{\fontencoding{T1}\fontfamily{#1}\selectfont}
%% if you don’t have the FontSite fonts either \renewcommand*{\FSfont}[1]{}
%% or use your own choice of family.
%% select a (TeX Font) font by its font family ID
\newcommand*{\TXfont}[1]{\fontencoding{T1}\fontfamily{#1}\selectfont}
%% Generic publisher’s logo
\newcommand*{\plogo}{\fbox{$\mathcal{PL}$}}
%% Some shades
\definecolor{Dark}{gray}{0.2}
\definecolor{MedDark}{gray}{0.4}
\definecolor{Medium}{gray}{0.6}
\definecolor{Light}{gray}{0.8}
%%%% Additional font series macros

\newcommand*{\titleLL}{\begingroup% Lost Languages
\drop=0.1\textheight
\fboxsep 0.5\baselineskip
\sffamily
\vspace*{\drop}
\centering
{\textcolor{Dark}{\HUGE Cinelerra-GG Version Infinity}}\par
\vspace{0.5\drop}
{\includegraphics[width=0.1\linewidth]{./images/cin-logo}}\par
\vspace{0.5\drop}
\colorbox{Dark}{\textcolor{white}{\normalfont\itshape\Large
User Manual}}\par
\vspace{0.3\drop}
{\Large Last update}\\
{\footnotesize \today}\\
\url{https://www.cinelerra-gg.org}\par
\vspace*{\drop}
{\includegraphics[width=0.5\linewidth]{./images/cin-big.png}}\par
\endgroup}

